The physical generation of lift is undoubtedly a viscous phenomenon, fully realized by the well known \textsc{Navier}--\textsc{Stokes} and continuity equations.
As the measure of selfadhesion, viscosity naturally induces vortical motion in the fluid.
Yet, a quasipotential formulation in which the effects of viscosity are realized, without the resulting motion being vortical, is what we shall utilize henceforth.
Assume then, that the fluid we investigate is inviscid, and that the fluid velocity is small such that the \textsc{Mach} number is small.
Upon assuming conservative body forces, the \textsc{Navier}--\textsc{Stokes} equation then reduces to the \textsc{Euler} and incompressibility equations\footnote{discuss \textsc{Mach} number, incompressibility is very relevant here.}
\begin{subequations}\label{eq:Euler_Incompressible}
  \begin{equation}\tag{\theequation{a, b}}
    \varrho\Dee_t \uvec = -\nabla p, \qquad \nabla \cdot \uvec = 0.
  \end{equation}
\end{subequations}
A vector field satisfying \subeqref{\ref{eq:Euler_Incompressible}}[b] is called solenoidal.
From the \textsc{Hodge} decomposition theorem, we may express the fluid velocity field as the sum $\uvec^{\flat} = \dee\Phi + \codee\vfrak$.
It is apparent that the curl of the \textsc{Euler} equation \subeqref{\ref{eq:Euler_Incompressible}}[a] is zero, such that the vector field describing the fluid flow is harmonic.
\begin{equation}\label{eq:Laplace_equation}
  \Delta \Phi = 0, \qquad \xvec \in \Omega
\end{equation}
This representation of the fluid is manageable and allows for powerful tools with which to investigate it, of which we shall promptly become aware.
For we shall no longer consider the differential equation above, but instead an integral formulation more fitting for computeraided solving.

Consider two harmonic scalar functions $\phi$ and $\varphi$.
Since all zero-forms are coclosed,%
\footnote{%
        The proof is trivial.
        $\codee\phi = \star\dee{\star}\phi = \dee\phi\wedge\dee V$, since the volume form is closed.
        The grade of $\codee\phi$ is then not congruent with the dimension of the manifold, and thus vanishes.
}
the \textsc{Hodge}--\textsc{Laplace} operator acting on zero-forms is equivalent to the familiar definition $\Delta\phi = \Div\Grad{\phi}$ of vector calculus.
Consider furthermore the expressions $\codee\phi\dee\varphi$ and $\codee\varphi\dee\phi$, the codifferential working on the product.
We take first the former, and expand the codifferential into its constituent stars and differential.
From the definition of the \textsc{Hodge} star, we know it to be impotent acting on a scalar field.
Thus,
\[
        \codee\phi\dee\varphi = \star \langle \dee \phi , \dee \varphi \rangle \,\dee V + \phi \Delta \varphi = \star \langle \dee \varphi , \dee \phi \rangle \,\dee V = \codee\varphi\dee\phi,
\]
following from the symmetry of the inner product.
Indeed, upon integrating their difference over the domain on which they are harmonic, we are free to exploit the fundamental theorem of multivariate calculus.
Remarkably, the behavior of harmonic functions is entirely determined by their boundary data.
This is in the special case of harmonic functions what is referred to as \textsc{Green}'s second identity, which in the notation of vector calculus, can be denoted
\begin{equation}\label{eq:Greens_second_identity}
  \int_{\partial\Omega} \big( \phi \nabla\varphi - \varphi \nabla\phi \big)\cdot\nhat \,\dee S = 0,
\end{equation}

\noindent
Constructing now a \textsc{Green} function defined to be harmonic everywhere, with the exception at some point $\bzhe \in \Othal$, we may recast equation \eqref{eq:Greens_second_identity} in terms of the \textsc{Cauchy} principal value.\footnote{explain}
For the case of three dimensional flow, the \textsc{Green} function is given by
\begin{subequations}\label{eq:Green_function_3D}
  \begin{equation}\tag{\theequation{a, b}}
    \green(\xvec; \bzhe) = \frac{1}{4\pi r}, \qquad r = \absl{\xvec - \bzhe}.
  \end{equation}
\end{subequations}
We confirm that $\Delta \green(\xvec; \bzhe) = \delta(\xvec - \bzhe)$, where $\delta$ is the \textsc{Dirac} delta function, meaning the \textsc{Green} function of equation \subeqref{\ref{eq:Green_function_3D}}[a] is harmonic almost everywhere.
We may insert this function for $\varphi$ in equation \eqref{eq:Greens_second_identity}, with the appropriate adjustments to the right-hand side.
By integrating around the singularity at $\bzhe$, we find that
\begin{equation}\label{eq:integral_equation}
  \principalvalue \int_{\partial\Omega} \big( \phi \nabla_{\nhat} \green - \green \nabla_{\nhat} \phi \big) \,\dee S = c\phi(\bzhe),
\end{equation}

\noindent
where $c$ is some constant dependent on the location of $\bzhe$.
For the directional derivative, we mean $\nabla_{\nhat} = \nhat \cdot \nabla$.
This is the basis for the boundary element method, the numerical implementation of which this thesis seeks to employ.
